\section{Transactions}
\label{sec:transactions}

A Bitcoin transaction consumes previously created outputs and creates new
outputs. Consensus validation distinguishes the serialized transaction bytes,
the identifiers derived from those bytes, and the contextual checks that depend
on the UTXO set and block height.

\subsection{Identifier and amount primitives}

\begin{longtable}{@{}L{0.22\linewidth}L{0.24\linewidth}L{0.42\linewidth}@{}}
\toprule
Type & Encoding & Meaning \\
\midrule
\code{Txid} & \code{uint256} & \(txid(tx)\), as defined in Section~\ref{sec:notation}. \\
\code{Wtxid} & \code{uint256} & \(wtxid(tx)\), as defined in Section~\ref{sec:notation}. \\
\code{OutPoint} & \code{Txid || uint32} & Reference to one transaction output by transaction identifier and output index. \\
\code{Amount} & \code{int64} & Satoshis. Consensus-valid output amounts satisfy $0 \le value \le$ \code{MAX\_MONEY}. \\
\bottomrule
\end{longtable}

\subsection{Transaction records}

\begin{longtable}{@{}L{0.18\linewidth}L{0.22\linewidth}L{0.22\linewidth}L{0.28\linewidth}@{}}
\toprule
\multicolumn{4}{@{}l}{\textbf{\code{TxOut}}} \\
\midrule
Field & Type & Encoding & Meaning \\
\midrule
\code{value} & \code{Amount} & \code{int64} & Output value in satoshis. \\
\code{scriptPubKey} & \code{Script} & \code{varbytes} & Locking script. \\
\bottomrule
\end{longtable}

\begin{longtable}{@{}L{0.18\linewidth}L{0.22\linewidth}L{0.22\linewidth}L{0.28\linewidth}@{}}
\toprule
\multicolumn{4}{@{}l}{\textbf{\code{TxIn}}} \\
\midrule
Field & Type & Encoding & Meaning \\
\midrule
\code{previousOutput} & \code{OutPoint} & 36 bytes & Output spent by this input, or the null outpoint for coinbase. \\
\code{scriptSig} & \code{Script} & \code{varbytes} & Unlocking script or coinbase data. \\
\code{sequence} & \code{uint32} & 4 little-endian bytes & Finality, replacement, and relative-lock field. \\
\bottomrule
\end{longtable}

\begin{longtable}{@{}L{0.18\linewidth}L{0.22\linewidth}L{0.22\linewidth}L{0.28\linewidth}@{}}
\toprule
\multicolumn{4}{@{}l}{\textbf{\code{Transaction}}} \\
\midrule
Field & Type & Encoding & Meaning \\
\midrule
\code{version} & \code{int32} & 4 little-endian bytes & Transaction version. Versions 1 and 2 are standard by default relay policy; consensus accepts the integer subject to contextual rules. \\
\code{inputs} & \code{vector<TxIn>} & CompactSize count, then inputs & Spent outputs or coinbase input. \\
\code{outputs} & \code{vector<TxOut>} & CompactSize count, then outputs & Created outputs. \\
\code{witnesses} & \code{vector<Witness>} & witness serialization only & One witness stack per input when witness serialization is used. \\
\code{lockTime} & \code{uint32} & 4 little-endian bytes & Absolute height or time lock. \\
\bottomrule
\end{longtable}

\code{Witness} is \code{vector<varbytes>}. It is present only in witness
serialization and has exactly one stack per input.

\subsection{Legacy serialization}

\(\mathrm{E}_{0}(tx)\), the legacy transaction serialization, is:

\[
  version || inputs || outputs || lockTime.
\]

\begin{longtable}{@{}L{0.22\linewidth}L{0.24\linewidth}L{0.42\linewidth}@{}}
\toprule
Part & Encoding & Meaning \\
\midrule
\code{version} & \code{int32} & Transaction version. \\
\code{inputCount} & \code{compactSize} & Number of inputs. \\
\code{inputs} & \code{TxIn[inputCount]} & Serialized inputs in order. \\
\code{outputCount} & \code{compactSize} & Number of outputs. \\
\code{outputs} & \code{TxOut[outputCount]} & Serialized outputs in order. \\
\code{lockTime} & \code{uint32} & Absolute lock field. \\
\bottomrule
\end{longtable}

\subsection{Witness serialization}

\(\mathrm{E}_{w}(tx)\), the witness transaction serialization, is:

\[
  version || marker || flag || inputs || outputs || witnesses || lockTime.
\]

\begin{longtable}{@{}L{0.22\linewidth}L{0.24\linewidth}L{0.42\linewidth}@{}}
\toprule
Part & Encoding & Meaning \\
\midrule
\code{marker} & \code{0x00} & Distinguishes witness serialization from legacy serialization. \\
\code{flag} & \code{0x01} & Indicates witness data follows. Other nonzero flag bits are invalid. \\
\code{witnesses} & \code{Witness[inputCount]} & One witness stack per input. \\
\bottomrule
\end{longtable}

The witness form is valid only when at least one witness stack is nonempty. A
transaction with no witness data uses legacy serialization for both \code{Txid}
and \code{Wtxid} \cite{bip-0141,bip-0144}.

\subsection{Transaction identifiers}

\begin{longtable}{@{}L{0.22\linewidth}L{0.58\linewidth}@{}}
\toprule
Identifier & Definition \\
\midrule
\code{txid(tx)} & Defined in Section~\ref{sec:notation}. \\
\code{wtxid(tx)} & Defined in Section~\ref{sec:notation}. \\
\code{GenTxid} & Pair of identifier kind and hash used by relay to distinguish txid and wtxid announcements. \\
\bottomrule
\end{longtable}

\subsection{Context-free validity}

A transaction is context-free valid only if:

\begin{enumerate}[leftmargin=*]
  \item it has at least one input and at least one output;
  \item its stripped size satisfies
        $strippedSize(tx) \cdot
        \text{\code{WITNESS\_SCALE\_FACTOR}}
        \le \text{\code{MAX\_BLOCK\_WEIGHT}}$;
  \item every output value is in range;
  \item the sum of output values is in range;
  \item no two inputs reference the same previous output;
  \item if it is coinbase, it has exactly one input whose previous output is
        the null outpoint and whose \code{scriptSig} length is between
        \code{COINBASE\_SCRIPTSIG\_MIN} and
        \code{COINBASE\_SCRIPTSIG\_MAX} bytes; and
  \item if it is not coinbase, no input uses the null outpoint.
\end{enumerate}

Context-free validity does not prove the inputs exist or signatures are valid;
those checks require the UTXO set and rule context
\cite{bitcoin-core-v31}.

\subsection{Coinbase transaction}

A coinbase transaction has exactly one input whose previous output is the null
outpoint defined by \code{NULL\_TXID} and \code{NULL\_INDEX}. Its input does not
spend a UTXO. Its outputs create the block subsidy and collect fees, subject to
the coinbase amount rule during block connection.

Coinbase-created outputs require \code{COINBASE\_MATURITY} confirmations before
they may be spent.

\subsection{Value conservation}

For a non-coinbase transaction:

\[
  fee(tx) = \sum value(spentInputs(tx)) - \sum value(outputs(tx)).
\]

The transaction is invalid if any referenced input is absent, any spent coin is
immature coinbase, any output is outside the money range, the output sum is
outside the money range, or \code{fee(tx)} is negative
\cite{bitcoin-core-v31}.

\subsection{Locktime and sequence}

\code{lockTime} values below \code{LOCKTIME\_THRESHOLD} are block heights;
values at or above that threshold are Unix times. For a candidate block at
height \(h\) with rule context \(R\) and locktime cutoff \(t\), define the
selected lock boundary:

\[
\ell(tx,h,t)=
\begin{cases}
h,
  & tx.\text{\code{lockTime}} < \text{\code{LOCKTIME\_THRESHOLD}},\\
t,
  & tx.\text{\code{lockTime}} \ge \text{\code{LOCKTIME\_THRESHOLD}},
\end{cases}
\]

\[
\begin{aligned}
\operatorname{AbsFinal}(tx,h,t) \Longleftrightarrow\;&
tx.\text{\code{lockTime}}=0 \vee
tx.\text{\code{lockTime}} < \ell(tx,h,t) \vee {}\\
&\forall i:\ tx.\text{\code{inputs}}[i].\text{\code{sequence}}
  = \text{\code{SEQUENCE\_FINAL}}.
\end{aligned}
\]

If \(R\) includes BIP113, \(t=\operatorname{MTP}(C,h-1)\); otherwise \(t\) is
the candidate block timestamp.

If \(R\) includes BIP68, relative locks apply only to version-2-or-higher
transactions and only to inputs whose \code{sequence} does not set
\code{SEQUENCE\_LOCKTIME\_DISABLE\_FLAG}. For validation view \(V\), define
\(\operatorname{SeqFinal}(tx,C,V,h)\) as follows. If
\(tx.\text{\code{version}} < 2\), it is true. Otherwise, for each applicable
input \(i\), let \(s_i\) be its sequence value, \(v_i=s_i \mathbin{\&}
\text{\code{SEQUENCE\_LOCKTIME\_MASK}}\), and \(c_i\) be the height of
the coin spent by input \(i\) in \(V\). Height locks require:

\[
  h > c_i + v_i - 1.
\]

When \(s_i\) sets \code{SEQUENCE\_LOCKTIME\_TYPE\_FLAG}, the input is time
locked instead. Let \(\tau_i=\operatorname{MTP}(C,c_i-1)\). Time locks require:

\[
  \operatorname{MTP}(C,h-1) >
  \tau_i + v_i \cdot \text{\code{SEQUENCE\_LOCKTIME\_GRANULARITY}} - 1.
\]

The transaction satisfies relative locks iff every applicable input satisfies
its selected height or time constraint; disabled inputs impose no constraint.
If \(R\) does not include BIP68, \(\operatorname{SeqFinal}\) is true
\cite{bip-0068,bip-0112,bip-0113,bitcoin-core-v31}.

\subsection{Signature hashes}

Signature validation signs a digest derived from the transaction, input index,
spent output, hash type, and active script version. Legacy, SegWit v0, and
Taproot/Tapscript spend paths use different digest algorithms, specified in
Section~\ref{sec:script}. These digests are consensus-critical because the
signature checks validate exactly those bytes.
